% SMA Crossover Backtest — research-note PDF.
%
% PROSE AND CODE ARE VERBATIM from the website article:
%   site/src/content/articles/sma-crossover-backtest.tsx
% The article interpolates its computed results into the prose; every
% interpolated value here is the one the site renders, read from
% data/sma-crossover-backtest.ts. Code cells are inline in the article.
%
% CHARTS generated by figures/make_sma_crossover_backtest.py from that module.
% The Sharpe grids carry a -9.99 sentinel where fast >= slow; those cells are
% masked and shown as an em dash rather than plotted.
\documentclass{pyportfolios-note}

\noteshorttitle{SMA Crossover Backtest}

\begin{document}
\thispagestyle{notecover}

\notemasthead{Quantitative Insights}{Tutorial}{Q3 2026}{}

\notetitle{SMA Crossover\\Backtest}

\notedeck{One rule, two parameters, a century of folklore — and every
backtesting sin visible on a single screen of code.}

\notelede{Buy when the 50-day average crosses above the 200-day; step aside when
it crosses back. That is the \textbf{golden cross} — the systematic
trend-following workhorse: one rule, two parameters, a century of folklore. Its
simplicity is exactly why it makes the perfect classroom for backtesting
discipline 101 — every sin that quietly ruins real backtests fits on one screen
of code here, where you can watch it happen. We run the same long-or-flat rule on
QQQ and BTC-USD over 2015–2024 with next-day execution and 10bp per-side costs,
then do what most crossover backtests skip: sweep the parameter grid, stress the
costs, and admit what one in-sample decade can and cannot prove.}

\notecard
  {Algorithmic Trading}
  {Pandas · NumPy · vectorbt}
  {QQQ + BTC-USD}
  {2015 – 2024}
  {Backtesting discipline — lookahead, costs, parameter sweeps and in-sample
   honesty}

\begin{multicols}{2}
\begin{notepipeline}
\item Download QQQ + BTC-USD closes, 2015–2024 (yfinance), dropna per asset
\item Signal: 50d SMA above 200d SMA → long, else flat
\item Execute next day — signal.shift(1), 10bp per side
\item Compare vs buy-and-hold; cross-check fills with vectorbt
\item Sweep the 4×4 fast/slow Sharpe grid per asset
\item Stress costs (0/10/25bp) and file the in-sample caveats
\end{notepipeline}
\end{multicols}

% ============================== 1 The rule: two moving averages =============
\noteneedroom{190bp}
\notesection{1.\notenumsep The rule: two moving averages}

\begin{multicols}{2}
Compute a fast and a slow simple moving average of the close. Think of the slow
one as the market's long memory and the fast one as its recent mood. When the
fast SMA sits above the slow one, recent prices sit above the long trend — be
long. When it drops below, be flat. No shorting, no leverage, no discretion, no
opinions. With \texttt{fast = 50} and \texttt{slow = 200} this is the golden
cross of financial television fame — and, more respectably, a one-asset special
case of \textbf{time-series momentum}, the effect Moskowitz, Ooi and Pedersen
documented across 58 futures markets.

We run it on two deliberately different animals: the Nasdaq-100 ETF (exchange
calendar, \textasciitilde252 bars a year) and Bitcoin (trades every day of the
week, \textasciitilde365 bars a year). BTC has prices on weekends where QQQ has
none, so each asset is backtested on its own calendar — per-asset
\texttt{dropna}, per-asset annualisation — never forward-filled onto a shared
grid.
\end{multicols}

\notecallout{Why practitioners care}{Trend filters run inside real money:
managed-futures funds, crypto treasuries, tactical ETF overlays. But the reason
this tutorial exists is the workflow, not the rule — the crossover is simple
enough that every backtesting sin (lookahead, ignored costs, parameter
cherry-picking, in-sample worship) is visible in a single screen of code. Learn
to catch them here and you will catch them in strategies that matter.}

% ================= 2 Next-day execution — the shift(1) that keeps you honest =
\noteneedroom{190bp}
\notesection{2.\notenumsep Next-day execution — the shift(1) that keeps you honest}

The most important line in the backtest is not the signal. It is the shift.
Today's SMA cross is computed on today's \textbf{close} — a price you cannot
trade before you have seen it. So today's signal may only earn
\textbf{tomorrow's} return. Drop the shift and the backtest buys every up-day one
bar early: a lookahead bias that flatters almost any signal and quietly
fabricates Sharpe.

\begin{notecode}
def backtest(px, fast=50, slow=200, fee=0.001):
    signal   = (px.rolling(fast).mean() > px.rolling(slow).mean()).astype(float)
    position = signal.shift(1).fillna(0.0)   # <- no lookahead. The whole game.
    ret      = px.pct_change().fillna(0.0)
    cost     = position.diff().abs().fillna(0.0) * fee   # 10bp per side
    return (position * ret - cost).iloc[250:]   # SMA warm-up window
\end{notecode}

Two more discipline details. Every position change pays 10bp — entries and exits
both. And the first 250 bars of each asset are reserved for SMA warm-up, so every
parameter pair we test later is judged on the \textbf{same} evaluation window:
2015-12-30 → 2024-12-31 for QQQ, 2015-09-08 → 2024-12-31 for BTC. The pandas core
above is deliberately transparent; we rebuild the identical trades with
vectorbt's \texttt{Portfolio.from\_signals} as an independent referee — final
value, max drawdown and trade count agree.

% ===================================== 3 Same rule, two verdicts ============
\noteneedroom{190bp}
\notesection{3.\notenumsep Same rule, two verdicts}

Growth of \$100, log scale (BTC would render every other line invisible
otherwise). Watch \textbf{where} the strategy detaches from buy-and-hold: always
in the big drawdowns. A slow trend filter is not an accelerator — it is a brake
pedal.

\notefigure{figures/sma-crossover-backtest/fig1-qqq-equity.png}
  {QQQ — growth of \$100 (log scale), 2015-12-30 → 2024-12-31}
  {The strategy tracks buy-and-hold closely and sidesteps part of the 2022
   drawdown}
  {Source: Author calculations, from the article's computed results.}

\notefigure{figures/sma-crossover-backtest/fig2-btc-equity.png}
  {BTC-USD — growth of \$100 (log scale), 2015-09-08 → 2024-12-31}
  {The strategy line steps flat through the 2018 and 2022 bear markets while
   buy-and-hold collapses and recovers}
  {Source: Author calculations, from the article's computed results.}

\begin{notetable}{@{}p{146bp}>{\raggedleft\arraybackslash}p{58bp}>{\raggedleft\arraybackslash}p{56bp}>{\raggedleft\arraybackslash}p{50bp}>{\raggedleft\arraybackslash}p{56bp}>{\raggedleft\arraybackslash}p{48bp}>{\raggedleft\arraybackslash}p{46.8bp}@{}}
\thcell{} & \thcell{Ann ret} & \thcell{Ann vol} & \thcell{Sharpe} & \thcell{Max DD} & \thcell{Trades} & \thcell{In mkt} \\
\theadrule
QQQ · 50/200 crossover & 17.9\% & 18.9\% & 0.97 & -28.6\% & 5 & 82\% \\ \rowrule
QQQ · buy \& hold & 19.0\% & 22.2\% & 0.90 & -35.1\% & — & 100\% \\ \rowrule
BTC · 50/200 crossover & 79.3\% & 59.3\% & 1.28 & -69.3\% & 9 & 64\% \\ \rowrule
BTC · buy \& hold & 89.6\% & 69.1\% & 1.27 & -83.4\% & — & 100\% \\
\end{notetable}

\begin{multicols}{2}
\begin{notebullets}
\item \textbf{QQQ} — lost the return race (\$439 vs \$478), but at lower vol and
      a third-shallower worst drawdown, for a slightly better Sharpe. No return
      edge; a brake pedal.
\item \textbf{BTC} — kept pace with one of the great bull markets in modern data
      while sitting out 36\% of all days, and cut the max drawdown from -83.4\%
      to -69.3\%.
\end{notebullets}
\end{multicols}

Neither run is a money machine; both are drawdown insurance. Two conservatisms
are baked into every number above: flat days are credited \textbf{nothing} (park
the cash at the T-bill rate and the strategy line improves while buy-and-hold's
cannot), and Sharpe is quoted on raw rather than excess returns. Both choices
shade against the strategy, which is the right direction to be wrong in.

% ============================ 4 Discipline 101: the parameter grid ==========
\noteneedroom{190bp}
\notesection{4.\notenumsep Discipline 101: the parameter grid}

One backtest is an anecdote. If 50/200 were genuinely special, small changes to
the parameters should not change the story — real effects are smooth
neighbourhoods, and only overfit ones are sharp peaks. So before believing
50/200, ask whether the \textbf{neighbourhood} agrees: we sweep fast ∈ \{10, 20,
50, 100\} × slow ∈ \{100, 150, 200, 250\} — 15 valid systems per asset,
identical execution, identical costs — and heat-map the Sharpe. Both grids share
one colour scale, so a glance tells you which asset rewarded trend.

\notefigure{figures/sma-crossover-backtest/fig3-qqq-grid.png}
  {QQQ — Sharpe by fast (rows) × slow (cols), 10bp per side · buy-and-hold = 0.90}
  {Every cell lands between 0.78 and 1.01, straddling buy-and-hold. Dashes mark
   pairs where fast ≥ slow}
  {Source: Author calculations, from the article's computed results. The rule on
   the colour bar marks the buy-and-hold Sharpe.}

\notefigure{figures/sma-crossover-backtest/fig4-btc-grid.png}
  {BTC-USD — Sharpe by fast (rows) × slow (cols), 10bp per side · buy-and-hold = 1.27}
  {The whole grid sits in a 1.18–1.43 band, mostly above buy-and-hold}
  {Source: Author calculations, from the article's computed results.}

This is the card's thesis in one picture. \textbf{QQQ}: every cell lands between
0.78 and 1.01, straddling buy-and-hold's 0.90 — only 7 of 15 cells beat it, none
decisively. Picking the best cell after the fact and calling it edge is selection
bias, not alpha; on this asset the honest claim is ``no return edge, a consistent
drawdown brake''. \textbf{BTC}: 12 of 15 cells sit above buy-and-hold's 1.27, the
whole grid stays in the 1.18–1.43 band, and every cell slashes the max drawdown.
The same rule gets two different verdicts because \textbf{trend behaves
differently per asset} — Bitcoin's decade delivered longer, cleaner trends and
far deeper crashes for the filter to sidestep.

% ======================================= 5 The costs dial: 0 / 10 / 25 bp ===
\noteneedroom{190bp}
\notesection{5.\notenumsep The costs dial: 0 / 10 / 25 bp}

Costs are the second-biggest backtest killer after lookahead — but they bite in
proportion to turnover. The 50/200 pair trades a handful of times a decade, so
even 25bp per side barely dents it. Speed the system up to 10/100 and the toll
booth opens:

\notefigure{figures/sma-crossover-backtest/fig5-costs.png}
  {Sharpe vs per-side cost, 0 → 25bp — slow pairs run flat, fast pairs pay the toll}
  {Both 50/200 lines are nearly horizontal; both 10/100 lines slope down as costs
   rise}
  {Source: Author calculations, from the article's computed results.}

\begin{notetable}{@{}p{164bp}>{\raggedleft\arraybackslash}p{78bp}>{\raggedleft\arraybackslash}p{88bp}>{\raggedleft\arraybackslash}p{78bp}>{\raggedleft\arraybackslash}p{76.8bp}@{}}
\thcell{} & \thcell{Trades} & \thcell{Sharpe @ 0bp} & \thcell{@ 10bp} & \thcell{@ 25bp} \\
\theadrule
QQQ · 50/200 & 5 & 0.97 & 0.97 & 0.96 \\ \rowrule
QQQ · 10/100 & 18 & 0.90 & 0.88 & 0.84 \\ \rowrule
BTC · 50/200 & 9 & 1.29 & 1.28 & 1.28 \\ \rowrule
BTC · 10/100 & 21 & 1.44 & 1.43 & 1.42 \\
\end{notetable}

Slow trend is nearly free to run. Fast trend pays a visible tax — QQQ's 10/100
variant loses 0.06 Sharpe going from free execution to 25bp, on 18 round-trip
entries. Any strategy whose backtest only works at 0bp does not work.

% ===================================== 6 One sample, no walk-forward ========
\noteneedroom{190bp}
\notesection{6.\notenumsep One sample, no walk-forward}

Everything above is one ten-year window, evaluated in-sample. There is no
walk-forward, no out-of-sample holdout, and no multiple-testing haircut for the
15 variants we just eyeballed per asset. 2015–2024 handed both assets two of the
strongest trend decades they have ever printed — a regime gift the next decade
owes nobody. Before promoting any cell of that grid, run the standard honesty
checklist: point-in-time data, walk-forward splits, pessimistic costs, and a
\textbf{deflated Sharpe} — Bailey \& López de Prado's correction for exactly the
selection bias a 15-cell grid search manufactures — on whatever looked best.

\notecallout{Practitioner take}{Trend filters are \textbf{regime insurance, not
alpha machines}. Priced honestly, the 50/200 bought you a third less drawdown on
QQQ at the cost of lagging a bull market, and on BTC it kept the upside while
skipping two 70\%+ crashes — that is an insurance payout, not stock-picking
genius. And note what actually moved the numbers in this tutorial: one
\texttt{shift(1)} and a costs assumption. Lookahead and ignored costs kill more
backtests than bad ideas ever will.}

\noteneedroom{190bp}
\begin{notereferences}
\item Brock, W., Lakonishok, J. \& LeBaron, B. (1992). Simple Technical Trading
      Rules and the Stochastic Properties of Stock Returns. \textit{Journal of
      Finance} 47(5).
\item Moskowitz, T., Ooi, Y.H. \& Pedersen, L.H. (2012). Time Series Momentum.
      \textit{Journal of Financial Economics} 104(2).
\item Bailey, D.H. \& López de Prado, M. (2014). The Deflated Sharpe Ratio:
      Correcting for Selection Bias, Backtest Overfitting and Non-Normality.
      \textit{Journal of Portfolio Management} 40(5).
\item vectorbt documentation — \texttt{Portfolio.from\_signals}, the
      signal-based backtesting API used as the fill cross-check.
\item Companion notebook: \texttt{sma-crossover-backtest.ipynb} — reproduces
      every figure from raw data (deterministic, no simulation).
\end{notereferences}

\end{document}
